\chapter{Desarrollo del Trabajo}\label{ch:desarrollo}

\section{Generación del Dataset Cuántico (Qiskit)}
\subsection{Definición de Circuitos Ansatz}
\subsection{Simulación de Operadores de Ruido}

\section{Arquitectura del Modelo de Aprendizaje (PyTorch)}
\subsection{Diseño de la Red}
% Trade-off entre capacidad del modelo (complejidad) y riesgo de sobreajuste.
\subsection{Función de Pérdida (Loss Function)}
% Justificación de MSE vs Fidelidad u otras distancias matriciales.

\section{Ciclo de Vida MLOps y Trazabilidad}
\subsection{Control de Semillas y Reproducibilidad}
\subsection{Tracking de Experimentos (MLflow/W\&B)}
% Análisis de las ejecuciones, Learning Rates y convergencia.

\section{Resultados y Evaluación Analítica}
\subsection{Rendimiento In-Distribution}
\subsection{Validación Out-Of-Distribution}
% Crucial: Demostrar que el modelo aprende la estructura del ruido y no memoriza los datos de entrenamiento.